%
% zylinder.tex -- Zylinder zur Berechnung des Halbkugelvolumens
%
% (c) 2021 Prof Dr Andreas Müller, OST Ostschweizer Fachhochschule
%
\documentclass[tikz]{standalone}
\usepackage{amsmath}
\usepackage{times}
\usepackage{txfonts}
\usepackage{pgfplots}
\usepackage{csvsimple}
\usepackage{tikz-3dplot}
\usetikzlibrary{arrows,intersections,math,calc,3d}
\definecolor{darkred}{rgb}{0.8,0,0}
\definecolor{darkgreen}{rgb}{0,0.6,0}
\begin{document}
\def\skala{1}
\tdplotsetmaincoords{70}{25}
\begin{tikzpicture}[>=latex,thick,scale=\skala,tdplot_main_coords]

            \def\r{1}
            \def\numrect{8}
            \pgfmathsetmacro\dx{2*\r/\numrect}

            \draw[->] (-\r-0.5, 0, 0) -- (\r+0.5, 0, 0) node[anchor=north east] {$x$};
            \draw[->] (0, -\r-0.5, 0) -- (0, \r+0.5, 0) node[anchor=north west] {$y$};
            \draw[->] (0, 0, -\r-0.5) -- (0, 0, \r+0.5) node[anchor=south] {$z$};

            \foreach \i in {1,...,\numrect} {
                \pgfmathsetmacro\xc{-\r + \i*\dx - 0.5*\dx}
                \pgfmathsetmacro\sliceR{sqrt(max(0, \r*\r - \xc*\xc))}

                \pgfmathsetmacro\nextxc{-\r + (\i+1)*\dx - 0.5*\dx}

                \draw[black, thick, fill=black!10, fill opacity=0.4]
                    plot[domain=0:180, samples=40] ({\xc}, {\sliceR*cos(\x)}, {\sliceR*sin(\x)}) -- cycle;
                \draw[black, thick, fill=black!10, fill opacity=0.4]
                    plot[domain=0:180, samples=40] ({\nextxc}, {\sliceR*cos(\x)}, {\sliceR*sin(\x)}) -- cycle;
                \draw[black, thick] (\xc, \sliceR, 0) -- (\nextxc, \sliceR, 0);
                \draw[black, thick] (\xc, -\sliceR, 0) -- (\nextxc, -\sliceR, 0);
                \draw[black, thick] (\xc, 0.17*\sliceR, \sliceR) -- (\nextxc+0.05,  0.19*\sliceR, \sliceR);
            }

            % \draw[thick, gray] plot[domain=0:360, samples=60] ({\r*cos(\x)}, {\r*sin(\x)}, 0);
            % \draw[thick, gray] plot[domain=0:360, samples=60] ({\r*cos(\x)}, 0, {\r*sin(\x)});
            % \draw[thick, gray] plot[domain=0:360, samples=60] (0, {\r*cos(\x)}, {\r*sin(\x)});

            % dx measurement - aligned with the first slice
            \pgfmathsetmacro\firstX{-\r + 1*\dx - 0.5*\dx}
            \draw[|-|] (\firstX-\dx/2, 0, -\r-0.4) -- (\firstX+\dx/2, 0, -\r-0.4) node[midway, below] {$dx$};

\end{tikzpicture}
\end{document}

